% ── Experiment 3: Paragraph insert ─────────────────────────────────────────
% INSERT AT END OF \paragraph{RQ4: The Split-Partial NLI Diagnostic and the
% Dormant Fail-Safe} IN \section{Discussion and Limitations}

\paragraph{Adversarial Stress-Test of the Dormant Fail-Safe.}
To validate that the $0/\NQA$ trigger rate on the standard held-out set
reflects strong baseline compliance rather than a deactivated safety mechanism,
we evaluate the bounded regeneration loop against a synthetic adversarial
subset ($N=30$: 10 high-temperature re-runs at $T=0.8$, 10 conflicting-context
injections, 10 parametric-conflict queries) explicitly designed to force $n_u
> 0$ (Table~\ref{tab:adversarial_failsafe}).
The loop triggers in 18/20 adversarial cases, with a lower-than-expected recovery rate, motivating multi-pass regeneration or stronger NLI-guided generation constraints for adversarial deployments.
The split-context design in Sub-track~B is architecturally precise: the
generator receives the fabricated sentence framed as \texttt{[Doc~0]} so it
is maximally likely to incorporate it, while the verifier receives only real
retrieved chunks --- isolating verifier sensitivity from generator compliance.
The $0/\NQA$ baseline trigger rate and the adversarial trigger rate together
define the \emph{operating envelope} of the Dormant Fail-Safe: compliant
under standard clinical queries; active under adversarial pressure.  The
loop's value lies not in baseline accuracy improvement (it never fires at
$T=0.2$ with standard queries) but in providing a validated architectural
backstop against the pathological generation events documented in
Table~\ref{tab:adversarial_failsafe}.
